% ---------------------------------------------------------------------------
\begin{frame}{What is a Functional Dependency?}
  \begin{block}{Definition}
    A \textbf{functional dependency (FD)} $A \to B$ holds in a relation $R$ if:
    \[
      \forall\, t_1, t_2 \in R:\quad t_1[A] = t_2[A] \;\Rightarrow\; t_1[B] = t_2[B]
    \]
    Whenever two tuples agree on $A$, they \emph{must} agree on $B$.\\
    We say: \textbf{``$A$ functionally determines $B$''}.
  \end{block}

  \begin{columns}[t]
    \begin{column}{0.5\textwidth}
      \textbf{Notation}
      \begin{itemize}
        \item $A \to B$ \quad --- ``$A$ determines $B$''
        \item $AB \to C$ \quad --- both $A$ and $B$ together determine $C$
        \item $A \twoheadrightarrow B$ \quad --- multivalued dependency (4NF)
      \end{itemize}
    \end{column}
    \begin{column}{0.45\textwidth}
      \textbf{Key Insight}\\
      FDs describe \emph{semantic} constraints of the real world.
      They are determined by the \textbf{meaning} of the data, not by one
      particular snapshot of values.
    \end{column}
  \end{columns}
\end{frame}

% ---------------------------------------------------------------------------
\begin{frame}{Functional Dependencies --- Examples from the Music Domain}
  Using our running example: an online music catalogue.

  \vspace{6pt}
  \begin{center}
    \small
    \begin{tabular}{ll}
      \toprule
      \textbf{Functional Dependency} & \textbf{Reading} \\
      \midrule
      $\texttt{ArtistID} \to \texttt{Name, Country, DebutYear}$ & Each artist ID has exactly one name/country/debut \\
      $\texttt{AlbumID} \to \texttt{Title, ArtistID, Genre, ReleaseYear}$ & Album ID determines all album-level attributes \\
      $\texttt{AlbumID, TrackNumber} \to \texttt{Title, Duration}$ & Track is uniquely identified by album + track no. \\
      \bottomrule
    \end{tabular}
  \end{center}

  \vspace{6pt}
  \begin{examples}{Concrete Data}
    \scriptsize
    \begin{tabular}{llll}
      \toprule
      \textbf{AlbumID} & \textbf{TrackNumber} & \textbf{Title} & \textbf{ArtistID} \\
      \midrule
      A100 & 1 & Anti-Hero     & ART1 \\
      A100 & 2 & Karma         & ART1 \\
      A101 & 1 & As It Was     & ART2 \\
      \bottomrule
    \end{tabular}\\[4pt]
    \normalsize
    $\texttt{AlbumID} \to \texttt{ArtistID}$ holds: both A100 rows give ART1.
  \end{examples}
\end{frame}

% ---------------------------------------------------------------------------
\begin{frame}{Types of FDs: Trivial, Partial, Transitive}
  \begin{block}{Trivial FD}
    $AB \to B$ --- the dependent is a \emph{subset} of the determinant.
    Always true by definition; carries no useful information; can be ignored in
    analysis.
  \end{block}

  \begin{block}{Partial FD --- Target of 2NF}
    $AB \to C$ holds, but also $A \to C$ alone holds.
    $B$ is \textbf{redundant} in the determinant.\\[2pt]
    \textit{Example:} Key is \texttt{(AlbumID, TrackNumber)}, but
    $\texttt{AlbumID} \to \texttt{Genre}$ holds alone.\\
    Genre depends only on the album, not the track --- \textbf{partial dependency}.
  \end{block}

  \begin{block}{Transitive FD --- Target of 3NF}
    $A \to B$ and $B \to C$ (with $B$ non-prime) implies $A \to C$
    \emph{indirectly}.\\[2pt]
    \textit{Example:}
    $\texttt{AlbumID} \to \texttt{ArtistID} \to \texttt{DebutYear}$\\
    Debut year depends on the artist, not directly on the album.
  \end{block}
\end{frame}

% ---------------------------------------------------------------------------
\begin{frame}{Types of FDs: Full and Multivalued}
  \begin{block}{Fully Functional FD}
    $AB \to C$ holds, and \emph{neither} $A \to C$ nor $B \to C$ alone holds.
    Both attributes in the determinant are genuinely required.\\[2pt]
    \textit{Example:} $\texttt{AlbumID, TrackNumber} \to \texttt{Title}$\\
    Neither the album ID alone nor the track number alone uniquely identifies a
    track title.
  \end{block}

  \begin{block}{Multivalued FD --- Target of 4NF}
    $A \twoheadrightarrow B$ --- knowing $A$ determines an entire \emph{set} of
    $B$ values, independently of other attributes.\\[2pt]
    \textit{Example:} A musician can have multiple instruments AND multiple
    languages spoken.  These are \textbf{independent} multi-valued facts.
  \end{block}

  \begin{alertblock}{Summary: Which NF Eliminates Which Dependency Type?}
    \small
    \begin{tabular}{ll}
      \textbf{2NF}  & Eliminates partial FDs \\
      \textbf{3NF}  & Eliminates transitive FDs \\
      \textbf{BCNF} & Every determinant must be a superkey \\
      \textbf{4NF}  & Eliminates multivalued FDs \\
    \end{tabular}
  \end{alertblock}
\end{frame}

% ---------------------------------------------------------------------------
\begin{frame}{[GRAPHIC PLACEHOLDER] FD Diagram for the Music Example}
  \begin{alertblock}{INSERT GRAPHIC HERE}
    Add a functional dependency diagram (attribute graph) showing:
    \begin{itemize}
      \item \texttt{AlbumID} $\to$ \texttt{Title}, \texttt{Genre},
            \texttt{ReleaseYear}, \texttt{ArtistID}
      \item \texttt{ArtistID} $\to$ \texttt{Name}, \texttt{DebutYear}
      \item \texttt{(AlbumID, TrackNumber)} $\to$ \texttt{Title}, \texttt{Duration}
      \item Dashed arrow from \texttt{AlbumID} through \texttt{ArtistID} to
            \texttt{DebutYear} showing the transitive chain
    \end{itemize}
    Suggested tool: TikZ (attribute nodes as rounded rectangles, solid arrows for
    direct FDs, dashed arrow for the transitive dependency)
  \end{alertblock}
\end{frame}
